% tex/preamble.tex
% Layout, Kästen und Makros für beide Teile des Dokuments.
% Wird von tex/main.tex eingebunden. Nie direkt kompilieren: immer
% python3 werkzeuge/build.py aus dem Repo-Wurzelverzeichnis.

\usepackage[T1]{fontenc}
\usepackage[utf8]{inputenc}
\usepackage[ngerman]{babel}
\usepackage{lmodern}
\usepackage{amsmath,amssymb,mathtools}
\usepackage{siunitx}
\usepackage{enumitem}
\usepackage{booktabs}
\usepackage{multirow}
\usepackage{graphicx}
\usepackage{float}
\usepackage{geometry}
\usepackage{xcolor}
\usepackage{fancyhdr}
\usepackage{titlesec}
\usepackage{caption}
\usepackage[most]{tcolorbox}

% Farben: ruhige Palette, Theorie kühl, Lösungen warm, Warnungen deutlich.
\definecolor{phblau}{HTML}{1F4E79}
\definecolor{phblauhell}{HTML}{EAF1F8}
\definecolor{phgruen}{HTML}{2E6F4E}
\definecolor{phgruenhell}{HTML}{EAF4EE}
\definecolor{phorange}{HTML}{9C4A11}
\definecolor{phorangehell}{HTML}{FDF1E7}
\definecolor{phgrau}{HTML}{5A5A5A}
\definecolor{phgrauhell}{HTML}{F2F2F2}

\usepackage[colorlinks=true,
            linkcolor=phblau,
            urlcolor=phblau,
            citecolor=phblau,
            bookmarksnumbered=true,
            bookmarksopen=true,
            pdftitle={Physik: Theorie und Lösungen},
            pdfsubject={Theorie und Lösungen zu den Übungsaufgaben}]{hyperref}
\usepackage{cleveref}

% Schalter für die PDF-Varianten; generated/modus.tex setzt die Werte.
\newif\ifzeigetheorie   \zeigetheorietrue
\newif\ifzeigeloesungen \zeigeloesungentrue

% ---------------------------------------------------------------- Seite

\geometry{a4paper, top=2.4cm, bottom=2.4cm, left=2.5cm, right=2.5cm,
          headheight=14pt}
\setlength{\parskip}{0.5em}
\setlength{\parindent}{0pt}
\linespread{1.05}

\pagestyle{fancy}
\fancyhf{}
\renewcommand{\headrulewidth}{0.4pt}
\fancyhead[L]{\small\color{phgrau}\nouppercase{\rightmark}}
\fancyhead[R]{\small\color{phgrau}%
  \ifzeigetheorie\ifzeigeloesungen Theorie und Lösungen\else Theorie\fi
  \else Lösungen\fi}
\fancyfoot[C]{\small\color{phgrau}\thepage}

\titleformat{\part}[display]
  {\normalfont\huge\bfseries\color{phblau}}{}{0pt}{\Huge}
\titlespacing*{\part}{0pt}{0pt}{24pt}

\titleformat{\section}
  {\normalfont\Large\bfseries\color{phblau}}{\thesection}{0.6em}{}
  [\vspace{-0.5em}{\color{phblau!35}\rule{\linewidth}{0.8pt}}]
\titlespacing*{\section}{0pt}{1.6em}{0.9em}

\titleformat{\subsection}
  {\normalfont\large\bfseries\color{phgrau}}{\thesubsection}{0.6em}{}
\titlespacing*{\subsection}{0pt}{1.1em}{0.5em}

\captionsetup{font=small,labelfont={bf,color=phgrau}}

\setlist[itemize]{leftmargin=1.2em,itemsep=0.15em,topsep=0.2em}
\setlist[enumerate]{leftmargin=1.5em,itemsep=0.15em,topsep=0.2em}

% ---------------------------------------------------------------- Zahlen

\sisetup{
  per-mode=symbol,
  detect-all,
  quantity-product = \,,
  range-phrase = { bis },
  list-final-separator = { und },
  separate-uncertainty = true,
  output-decimal-marker = {,},
}

% ---------------------------------------------------------------- Makros

% Ableitungen und Differentiale: \dd{x}, \dv{x}{t}, \ddv{2}{x}{t}, \pdv{f}{x}
\newcommand*{\dd}[1]{\mathop{}\!\mathrm{d}#1}
\newcommand*{\dv}[2]{\frac{\dd{#1}}{\dd{#2}}}
\newcommand*{\ddv}[3]{\frac{\dd^{#1}{#2}}{\dd{#3}^{#1}}}
\newcommand*{\pdv}[2]{\frac{\partial {#1}}{\partial {#2}}}
\DeclarePairedDelimiter{\abs}{\lvert}{\rvert}
\DeclarePairedDelimiter{\norm}{\lVert}{\rVert}
\DeclarePairedDelimiter{\betrag}{\lvert}{\rvert}
\newcommand*{\vektor}[1]{\vec{#1}}
\newcommand*{\mittel}[1]{\overline{#1}}
\newcommand*{\Hin}{\text{Hinweis}}
\newcommand*{\konst}{\text{const.}}

% ---------------------------------------------------------------- Kästen

\tcbset{
  rahmen/.style={
    enhanced, breakable,
    boxrule=0.6pt,
    arc=2pt,
    left=8pt, right=8pt, top=7pt, bottom=7pt,
    fonttitle=\bfseries\small,
    before skip=0.6em, after skip=0.6em,
  },
}

% Werkzeugkasten: das, was man zum Selberlösen braucht, in Teil I.
\newtcolorbox{werkzeugkasten}[1][Werkzeuge]{rahmen,
  colback=phblauhell, colframe=phblau, coltitle=white,
  title={#1}}

% Typische Falle, Vorzeichen- und Einheitenfehler.
\newtcolorbox{achtung}[1][Typische Falle]{rahmen,
  colback=phorangehell, colframe=phorange, coltitle=white,
  title=#1}

% Ergebnis einer Rechnung, hervorgehoben.
\newtcolorbox{ergebnis}[1][Ergebnis]{rahmen,
  colback=phgruenhell, colframe=phgruen, coltitle=white,
  title=#1}

% Kontrolle: Einheit, Grenzfall, Größenordnung.
\newtcolorbox{kontrolle}[1][Kontrolle]{rahmen,
  colback=phgrauhell, colframe=phgrau, coltitle=white,
  title=#1}

% Merksatz zum Übertragen auf andere Aufgaben.
\newtcolorbox{merksatz}[1][Merksatz]{rahmen,
  colback=phgruenhell, colframe=phgruen!70!black, coltitle=white,
  title=#1}

% -------------------------------------------------- Kopf und Querverweise

% Kopfzeile einer Aufgabe in Teil I: Fundstelle im Skript.
\newcommand{\aufgabenkopf}[1]{%
  \par\vspace{-0.6em}\noindent{\small\color{phgrau}Skript: #1}\par\vspace{0.4em}}

% Verweis von der Theorie zur Lösung derselben Aufgabe.
\newcommand{\loesungsverweis}[1]{%
  \ifzeigeloesungen
    \par\vspace{0.4em}\noindent{\small\color{phblau}%
    \hyperref[loes:#1]{$\rightarrow$ Lösung dieser Aufgabe}}%
    \par\vspace{0.2em}%
  \fi}

% Verweis von der Lösung zurück in den Theorieteil.
\newcommand{\theorieverweis}[1]{%
  \ifzeigetheorie
    \par\vspace{0.4em}\noindent{\small\color{phblau}%
    \hyperref[aufg:#1]{$\leftarrow$ Werkzeuge dazu in Teil I}}%
    \par\vspace{0.2em}%
  \fi}

% Rückfallwerte, die generated/meta.tex überschreibt.
\providecommand{\SkriptStand}{Skript noch nicht hinterlegt}
\providecommand{\SkriptQuelle}{}
\providecommand{\StandDatum}{\today}
\providecommand{\Fachtitel}{Fach}
